\section{A Contribution to Understanding Prof. Georg Sverdrup and His Selected Writings}

\begin{center}
\begin{minipage}{0.7\textwidth}
\centering
When all that you meant was met with mockery,\\
and all that you wrought was only despised,\\
when your whole intent, however good it was,\\
is dismissed as nonsense by a self-satisfied fool,\\
then it must be practiced,\\
then it shall be tried,\\
your love.\\[0.5em]
Bjørnstjerne Bjørnson.
\end{minipage}
\end{center}

The little verse above, by our great poet, is found unsigned in the Collin manuscript collection. It was printed for the first time in 1907, the year of Professor Sverdrup’s death. We do not know when the poem was composed, or whether the poet wrote it for any particular person. We do know that scarcely anything could be more fitting as a tribute to Sverdrup, if we wish to express how great his work was, and how little it was understood.

A longer time must pass before history can give a many-sided assessment of Sverdrup’s life and work. As things now stand, there are far too many who have no understanding of what he accomplished among us. On the other side are those who seem to forget that Sverdrup too was subject to human frailties, and that now and then he was mistaken.

But there ought not to be divided opinions as to whether the goal at which Sverdrup aimed was the right one, and whether the principles upon which he built were sound. Likewise, no one should doubt that his whole labor was borne by earnestness, honesty, and love.

Although much has been said about the goal and about the principles, it is a fact that, to a great extent, they are both unknown and misunderstood. Among many of our people there has been a lack either of good will or of ability to concern themselves with them. Many have done what they could to lay obstacles in the way of the principles’ realization. Some have practiced them in such a manner that they have become mere caricatures. Others, demonstrative riders of principle, have degraded them into meaningless slogans. The well-deserved mockery has not failed to appear, though it has fallen more upon the innocent than upon the guilty.

Yet amid all this there are many who have labored with sobriety for, and in faithfulness to, the principles upon which Sverdrup laid such great weight: the doctrine of the child (Barnelovdommen), the congregational autonomy, and purposeful pastoral training. These principles existed both in theory and in practice before Augsburg came to be; but for the Norwegian people here as well as in Norway they were unknown. Augsburg joined them together and made them, already in the seventies, its sine qua non; the question of originality and priority we may leave out of account. What is of consequence is that Augsburg has done incomparable pioneer work in the Norwegian-American Church, and that it has impressed upon the consciousness of our people these truths: that life is more than doctrine, the congregation is more than a mere institution, pastor more than teacher and liturgist. And of consequence for Augsburg is this, that a theologian such as Sverdrup laid down his life in its service, thirty-three years of hard, self-denying, almost incomprehensible labor, which has borne and will bear blessed fruit for the Kingdom of God.

In Sverdrup’s Collected Writings there is now given a new occasion to make oneself acquainted with what Sverdrup willed, worked for, and meant. By reading his contributions and treatises, which extend over a period of thirty-eight years, one will see whether there was any breach in his development, any groping in his work, any self-contradiction in his principles. The work is planned for five or six volumes; it is therefore the most significant literary undertaking yet produced by Norwegian-American publishing.

Sverdrup has made so great a contribution to our ecclesiastical development that our Norwegian-American church people, in order to understand the present alignments among our church bodies, and in order better to find a remedy for certain ecclesiastical defects, ought and must make themselves familiar with the most important treatises from Sverdrup’s pen. Difference in church view and theological direction ought not to keep anyone from concerning himself with his collected writings, which concern the laity no less than the learned, and are no less intelligible to the former than to the latter. For everything Sverdrup wrote was clear, logical, and written for ordinary people.

His writings may be read with profit by all. But with special profit they will be read by his students, who heard him at the lectern.

In judging what a man has written, there are not a few who take as their standard the impression they have received of him as a speaker. Everything they must measure by the yardstick of eloquence. Such people, if for example they have heard Bjørnson and Ibsen speak, will naturally think more highly of Bjørnson’s literary productions than of Ibsen’s; for Bjørnson was as powerful in speech as he was with the pen, whereas Ibsen, the greatest dramatist of our age, was a poor speaker. A mixed audience will always produce rather different judgments. Most who read Otto Funcke’s writings will probably think that Funcke was a mighty speaker, which is not at all the case. And many have the notion that it was an imposing bearing, a magnetic personality, and brilliant gifts of speech that proved decisive in Paul’s great work; for Paul was indeed one of the greatest organizers known to world history. But we learn that his critics in Corinth maintained that his letters were weighty and strong, but his bodily presence weak and his speech to be set at naught (2 Cor. 10:10). He did not impress by Greek rhetoric, but he possessed profound Christian insight (2 Cor. 11:6).

Neither did Sverdrup impress by his outward appearance, or by “American Oratory”: there are those who find no trace of authority either in what he wrote or in what he spoke.

What was great in Sverdrup was his deep Christian insight. But he also knew how to present it with weight and force in speech. But where? Was it in the pulpit, on the platform, or at the lectern? 

Where did his strength as a speaker lie?

Before we answer this, we may mention a couple of interesting things that have been said about him and will be said about him still; for the wisdom of the street will never die.

A workingman once told me that he had once heard Sverdrup in the pulpit. “It was a dry sermon,” he said. A merchant received a pamphlet which Sverdrup had written about Augsburg’s principles. In irritation he burst out: “They do not go deep.” A pastor told me that he thought it strange that Sverdrup had so many admirers; for everything he wrote was poor as well. But—to return to our question.

It was not in the pulpit that Sverdrup’s proper strength lay. And yet he was a great preacher. When he preached, one received a sermon pervaded by the biblical spirit; it was a many-sided exposition and application of the text, fresh, original—Sverdrup never traveled the bare country road in his treatment; it had not a little in common with Old Testament prophetic speech. He did not stand in the pulpit in order to entertain, but to edify. He therefore preached without a trace of declamation, priestly unction, theatrical pathos—those artificial additives which, for many, are the essential thing in a sermon. One had to become accustomed to Sverdrup’s preaching. The more one heard him, the more one wished to hear him. For he opened the Scriptures.

Nor did his strong side appear upon the platform. His occasional speeches were free of theatrics, even in tone, unforced, and weighty with thought. They flowed easily, and the right word was never lacking; but it was too much marked by earnestness to serve up colorful fireworks, sparkling witticisms, brilliant flashes of wit.

If, on the other hand, when it came to debating a matter where he had opposition, then one learned to know him from a new side. Here one saw him as the ready and fearless disputant, who struck his opponents quickly, hard, and surely. Nothing could throw him off balance. He mowed down one row of arguments after another, and that with such self-command and force of conviction that his opponents could stare at him in mute, helpless bitterness. The annual meeting of the United Church, which in the nineties was given over almost entirely to debate and business, more than once showed that Sverdrup had been endowed from the cradle with something that had a disturbing effect on some of those who occupied the seats of honor in the assembly. To end the connection with Sverdrup, all constitutional procedure was set aside.

Sverdrup’s strongest side was the lectern, especially systematic theology. His other principal field was introductory Old Testament study and exegesis. Characteristic of his teaching in the Old Testament was his pragmatic presentation of the Lord’s dealings with Israel’s congregation, the gripping psychological portrayals of the Lord’s chosen instruments, the people’s pride, unbelief, chastisement, deliverance. In the Old Testament he was at home as few were, for he had penetrated into its spirit. A great deal of his work on the Old Testament survives, though it has never been printed. It would be a great loss if it were not included in the planned work. We mean especially his introduction and his commentary on the Psalms and the Prophets. Yet it was in systematic theology that Sverdrup, so to speak, took things by storm. For here he had use for his dialectic. It was especially with a view to his inspiring lectures in dogmatics that I wrote in an English paper in Philadelphia (1907) the following words—a small excerpt from a longer article:

\begin{quote}
His lectures captivated friend and foe alike by their brilliancy, force, and logic.... All was life, energy, enthusiasm when he was at the desk; his students were then all eye, all ear, exerting themselves to the utmost to get a verbatim report. Even the notes of his mediocre hearers are readable literature. His command of language was wonderful. He had the happy faculty of saying things in such a way that they were not to be forgotten. Every lecture was an oration, natural and unstudied, full of vim and vigor, with a most happy interplay of color and rhythm. Every sentence seemed to bristle with suggestions, every utterance breathed the spirit of liberty and convincing calmness.
\end{quote}
